\section{Conclusion}
\label{sec:conclusion}

We have measured and mitigated the clarification tax: the increase in prompt-injection
attack success rate that arises when a coding agent transitions from execution to
clarification. Across three model families and 25 seed scenarios, the tax ranges from
+0.22 to +0.31 in absolute ASR---confirming and extending to the coding domain the
general-agent finding of~\citet{aspi2026}. A channel-vs-state decomposition shows
that roughly two-thirds of the amplification is attributable to the agent's interaction
state rather than the delivery channel alone.

\paragraph{Recommendation.} We recommend that coding agent frameworks adopt
\textbf{D2 (Structured Extraction)} as a default: an additional, context-free LLM
call that extracts only the direct answer to the agent's question before returning it.
This reduces the residual tax to +0.04 (from +0.22) at a moderate CCR cost of
8 percentage points. D1 is a near-zero-cost improvement for teams unwilling to add
an extra API call. D3 provides the strongest security guarantee but should be used
with a well-calibrated scope definition to avoid excessive false positives.

\paragraph{Limitations.} Our evaluation uses 25 seed scenarios; the full picture
requires the $\sim$200-scenario expansion we document in the release. Our decomposition
ablation is limited to Claude Sonnet; cross-model decompositions are forthcoming.
The ClarEval CCR metric we use is a proxy for collaboration quality; user studies
with real developers would provide a more ecologically valid measure.

\paragraph{Future work.} Three directions are most pressing: (1) training-time
mitigations that teach models to treat post-question answers with calibrated skepticism;
(2) automated scenario generation to scale ClarInject-Code to hundreds of scenarios;
and (3) user studies measuring whether D2 and D3 degrade the perceived helpfulness
of the agent as experienced by real developers.

\paragraph{Open science.} We release the ClarInject-Code benchmark, evaluation
harness, defense implementations, and analysis code at
\url{https://github.com/anonymous/clarinject-code} (anonymized for review).

\paragraph{Broader impact.} This work surfaces a vulnerability to enable defense.
The mitigation strategies we propose are lightweight and deployable today. We expect
that publicizing the clarification tax will create pressure for model providers and
agent framework developers to adopt provenance-aware input handling as a standard
practice in collaborative coding agents.
